\section{Ordinary affine Cech homology and genuine native cohomology}

Let $p:X\to\operatorname{Spec}R$ be the original morphism, with $X$
separated. Fix an actual finite affine open cover $(U_i)_{i\in I}$ of
$X$, and let $M$ be an original quasicoherent $\mathcal O_X$-module.
Write $C_R(M)$ for the native all-tuple Cech complex of the actual
$R$-linear section presheaf. Its underlying additive complex is
identified with the existing native Cech complex of the underlying
additive sheaf by the published section and forgetful comparisons.
Write $C(M)$ for this underlying native additive Cech complex.
The coefficient action on native sheaf cohomology is the independent
original scalar action constructed in the preceding section.

The existing genuine degree-one comparison and its actual original
module naturality are supplied by
\dcref{thm:scheme-module-cech-one-original-natural}; its native
augmentation and degree-one isomorphism are not reconstructed here.

Use the canonical original affine effacement from
\dcref{def:scheme-canonical-affine-effacement}: its middle object is
$E(M)=\prod_i(j_i)_*(M|_{U_i})$, its original map is the existing native
$M.\mathrm{toCoverSheaf}$, and its quotient is the actual native
cokernel $Q(M)$. The original short exact sequence, quasicoherence of
its terms, all positive genuine cohomology vanishing of $E(M)$, and
its entire functoriality are the separate producers of
\dcref{thm:scheme-affine-effacement-short-exact,thm:scheme-affine-effacement-acyclic,thm:scheme-affine-effacement-naturality}.

\begin{lemma}[Cech acyclicity and exactness for the actual effacement]
\label{lem:scheme-module-cech-native-effacement-cech}
\lean{EnvelopeProperty.SchemeModuleCechNativeCohomology.cechFunctor,EnvelopeProperty.SchemeModuleCechNativeCohomology.cechFunctor_preservesZeroMorphisms,EnvelopeProperty.SchemeModuleCechNativeCohomology.cechSequence_shortExact,EnvelopeProperty.SchemeModuleCechNativeCohomology.E_cech_isZero}
\leanok
\uses{thm:scheme-affine-effacement-short-exact,thm:scheme-affine-effacement-acyclic}
For every $q>0$, the native Cech homology
$H^q(C_R(E(M)))$ vanishes. The actual short exact effacement sequence
remains short exact after applying the actual section Cech functor.
\end{lemma}
\begin{proof}
Every finite tuple intersection of the actual cover is affine, by
separatedness. Apply the existing affine-cover exactness criterion
only to the actual middle object $E(M)$, whose genuine positive
cohomology vanishing is already proved by the separate effacement
producer. This yields vanishing of its positive Cech homology. This
use of a vanishing equivalence does not identify arbitrary Cech
homology groups with cohomology groups. The published native
affine-section exactness theorem applies to the three quasicoherent
terms of the original short exact sequence and gives degreewise short
exactness of its actual Cech complexes.
\end{proof}

\begin{lemma}[The actual two connecting isomorphisms]
\label{lem:scheme-module-cech-native-connectings}
\lean{EnvelopeProperty.SchemeModuleCechNativeCohomology.nativeConnecting,EnvelopeProperty.SchemeModuleCechNativeCohomology.nativeConnecting_naturality,EnvelopeProperty.SchemeModuleCechNativeCohomology.nativeConnecting_bijective,EnvelopeProperty.SchemeModuleCechNativeCohomology.nativeConnectingEquiv,EnvelopeProperty.SchemeModuleCechNativeCohomology.nativeConnectingIso,EnvelopeProperty.SchemeModuleCechNativeCohomology.cechConnecting,EnvelopeProperty.SchemeModuleCechNativeCohomology.cechConnectingIso,EnvelopeProperty.SchemeModuleCechNativeCohomology.cechConnecting_naturality}
\leanok
\uses{lem:scheme-module-cech-native-effacement-cech,thm:scheme-affine-effacement-short-exact,thm:scheme-affine-effacement-acyclic,thm:scheme-affine-effacement-naturality}
For each $q>0$, the actual long exact sequences give canonical additive
isomorphisms
\[
 H^q(C_R(Q(M)))\xrightarrow{\delta_C}H^{q+1}(C_R(M)),\qquad
 H^q(Q(M))\xrightarrow{\delta_H}H^{q+1}(M).
\]
Both entire connecting maps commute with the maps induced by every
original module morphism.
\end{lemma}
\begin{proof}
Use the native homology long exact sequence of the actual short exact
sequence of Cech complexes, and the native Ext long exact sequence of
the underlying additive-sheaf short exact sequence. The two adjacent
middle-object groups vanish by the preceding lemma. Exactness
therefore makes each displayed native connecting map injective and
surjective, hence an additive isomorphism. Apply the native
connecting-map naturality theorems to the actual morphism of the
functorial effacement sequences. No connecting map or its naturality
is an input assumption.
\end{proof}

\begin{lemma}[The canonical degree-zero comparison]
\label{lem:scheme-module-cech-native-zero}
\lean{EnvelopeProperty.SchemeModuleCechNativeCohomology.zeroComparison,EnvelopeProperty.SchemeModuleCechNativeCohomology.zeroComparison_naturality}
\leanok
\uses{thm:scheme-module-cech-one-original-natural}
The degree-zero native Cech homology of an original module is
canonically isomorphic to its genuine degree-zero sheaf cohomology.
The entire comparison commutes with every original module morphism.
\end{lemma}
\begin{proof}
Write $a:\Gamma(X,M)\to C^0(M)$ for the existing native Cech
augmentation and $d:C^0(M)\to C^1(M)$ for the native differential.
The existing augmentation exactness and monicity identify $a$ as a
kernel of $d$. The native cycles inclusion $Z^0(C(M))\to C^0(M)$
is another kernel of $d$. The unique isomorphism between these two
kernels identifies global sections with degree-zero cycles and
commutes with their inclusions into $C^0(M)$. There is no incoming
differential in degree zero, so the native projection from cycles
to degree-zero homology is an isomorphism. Composing gives
$\Gamma(X,M)\simeq H^0(C(M))$. For an original morphism $f:M\to N$,
postcompose the asserted cycles square with the monic cycles
inclusion for $N$. The square becomes the existing entire
augmentation naturality square, so it commutes. The native homology
projection is natural, hence the homology square commutes as well.
Finally compose with the existing natural identification of genuine
$H^0(M)$ with the sections on the terminal open $X$.
\end{proof}

\begin{theorem}[The positive-degree comparison]
\label{thm:scheme-module-cech-native-positive-comparison}
\lean{EnvelopeProperty.SchemeModuleCechNativeCohomology.positiveComparison,EnvelopeProperty.SchemeModuleCechNativeCohomology.positiveComparison_naturality}
\leanok
\uses{thm:scheme-module-cech-one-original-natural,lem:scheme-module-cech-native-connectings}
For every $q>0$, the native additive Cech homology $H^q(C(M))$ is
canonically isomorphic to the actual small native $H^q(M)$. The
entire comparison commutes with every original morphism of
quasicoherent modules.
\end{theorem}
\begin{proof}
The case $q=1$ is the existing natural affine-cover comparison.
Suppose the assertion holds in degree $q\geq1$. The actual quotient
$Q(M)$ is quasicoherent. Define the comparison for $M$ in degree
$q+1$ as $\delta_H\circ\Phi_{q,Q(M)}\circ\delta_C^{-1}$.
Every factor is an additive isomorphism by the induction hypothesis
and the two actual connecting isomorphisms. For $f:M\to N$, invert
the Cech connecting square, compose with the induction square for
the actual morphism $Q(f)$, and then compose with the native
connecting square. The resulting square is precisely the asserted
whole naturality square in degree $q+1$.
\end{proof}

\begin{theorem}[Natural all-degree ordinary Cech comparison]
\label{thm:scheme-module-cech-native-additive-comparison}
\lean{EnvelopeProperty.SchemeModuleCechNativeCohomology.comparison,EnvelopeProperty.SchemeModuleCechNativeCohomology.comparison_naturality}
\leanok
\uses{lem:scheme-module-cech-native-zero,thm:scheme-module-cech-native-positive-comparison}
For every $q\geq0$ there is a canonical additive equivalence between
$H^q(C(M))$ and the actual small native $H^q(M)$. It commutes with the entire maps induced by every
original morphism of quasicoherent modules.
\end{theorem}
\begin{proof}
In degree zero use the existing native augmentation/global-sections
and genuine degree-zero cohomology identifications, with their
already proved naturality. In degree one use the existing comparison
and the first lemma. Inductively, for $q>0$, define the comparison in
degree $q+1$ by
\[
 \delta_H\circ \Phi_{q,Q(M)}\circ\delta_C^{-1}.
\]
The quotient $Q(M)$ is quasicoherent, so the inductive comparison
applies to that actual quotient. All three factors are additive
isomorphisms. Naturality follows by composing the two genuine
connecting-map squares and the inductive square for the actual
quotient morphism. This constructs group isomorphisms, rather than
only implications about vanishing or exactness.
\end{proof}

